% ============================================
% DOC-10 (ZH)
% 中文翻译 — SPEC-935-R356
% ============================================
\section*{\textbf{DOC-10} --- ECT记录}
\markboth{DOC-10 --- ECT记录}{}

\noindent \textit{37次无麻醉的治疗}

\noindent \textbf{巴尔巴塞纳收容医院 —— 激进治疗病区}
\textbf{患者:} 若阿金·安东尼奥·科雷亚 —— 登记号1,260
\textbf{期间:} 1937年至1952年

\vspace{0.5em}\noindent\rule{\textwidth}{0.5pt}\vspace{0.5em}

\subsection*{摘要}

\noindent 登记治疗次数: \textbf{37}
年平均: \textbf{2.46}
设备: Konvulsator IV型电休克机（德国，1932年）
麻醉: \textbf{未使用}（37次均未使用）
肌松药: \textbf{未使用}（37次均未使用）
事前镇静: \textbf{无}（患者被标注为“对言语镇静难治”）

\vspace{0.5em}\noindent\rule{\textwidth}{0.5pt}\vspace{0.5em}

\subsection*{详细记录（部分）}

\begin{moltablesix}
\hline
\# & 日期 & 电压 & 时长 & 即时效应 & 迟发效应 \\
\hline
1 & 37年3月12日 & 110V & 2.5秒 & 强直阵挛发作、口吐白沫、意识丧失12分钟 & 舌裂伤、尿失禁、意识混乱6小时 \\
2 & 37年3月19日 & 110V & 3.0秒 & 发作、T7椎体骨折（未治疗） & 慢性腰痛，患者持续陈述40余年 \\
3 & 37年3月26日 & 130V & 3.0秒 & 持续发作（4分钟）、发绀 & 48小时逆行性遗忘 \\
4 & 37年4月2日 & 130V & 3.5秒 & 发作、舌头严重咬伤 & 进食困难10天 \\
5 & 37年4月9日 & 130V & 3.5秒 & 发作、下颌脱位 & 无麻醉手法复位 \\
--- & --- & --- & --- & --- & --- \\
18 & 40年6月15日 & 150V & 4.0秒 & 发作、呼吸暂停（已复苏） & 患者陈述“飘出体外” \\
--- & --- & --- & --- & --- & --- \\
27 & 45年11月3日 & 150V & 4.0秒 & 发作、右股骨骨折 & 永久性跛行 \\
--- & --- & --- & --- & --- & --- \\
37 & 52年2月18日 & 170V & 5.0秒 & 最大发作、心跳骤停（已复苏） & 患者昏迷3天。醒来时说：“若阿金不在这里” \\
\hline
\end{moltablesix}

\vspace{0.5em}\noindent\rule{\textwidth}{0.5pt}\vspace{0.5em}

\subsection*{临床意见（转录）}

\noindent \textbf{内尔松·丰特斯医生，1937年:}
\textit{“患者对治疗无反应。妄想持续不变。将电压提高到130V。”}

\noindent \textbf{阿米尔卡·巴尔博萨医生，1940年:}
\textit{“18次治疗后，患者开始陈述电击期间的‘离体体验’。我认为这是一个进步——意识正在与身体脱离。继续。”}

\noindent \textbf{阿米尔卡·巴尔博萨医生，1945年:}
\textit{“治疗中出现股骨骨折。患者未感到骨折——只陈述‘听见骨头断了但不疼’。建议提高电压。患者的耐受性非同寻常。”}

\noindent \textbf{埃托尔·奥利维拉医生，1979年（30年后的页边注）:}
\textit{“这份病历是一份酷刑文件。37次无麻醉。骨折未治疗。没有知情同意。而患者扛过了这一切——清醒到足以写下六十二年的日记。这不是难治性。这是韧性。或者是附身。或者两者都是。”}

\vspace{0.5em}\noindent\rule{\textwidth}{0.5pt}\vspace{0.5em}

\subsection*{档案员注（2026）}

\noindent Konvulsator IV型设备的设计用途是配合麻醉和肌松药使用。在收容院，两者皆无。制造商推荐的最大电压是130V。患者1,260在最后一次治疗中接受了高达170V的电压。

\noindent 患者在没有明显脑损伤的情况下存活，这“在临床上无法解释”（托雷斯医生，2023年）。

\noindent 患者1,260在他的日记（1952年）中写道：\textit{“电击影响不了怪物。电击烧掉我的碎片。怪物吃掉剩下的。”}

\noindent 有人建议，患者对治疗“无法解释的”耐受性可能与其妄想内容有关——具体地说，与他对一个吸收创伤的内在实体的信念有关。

\noindent 或者，换一种说法，那个妄想并不是妄想。

\noindent 档案馆对后一种假设不持立场。

\vspace{0.5em}
\noindent\textit{\textbf{↪ 路线:} \rota{MEM-17} \rotasep \rota{DOC-02} \rotasep \rota{EPI-01}}
